\documentclass[11pt]{article}

\usepackage[margin=1in]{geometry}
\usepackage{amsmath,amssymb,amsthm}
\usepackage{graphicx}
\usepackage[hidelinks]{hyperref}

\newtheorem{theorem}{Theorem}
\newtheorem{corollary}{Corollary}

\title{A two-point counterexample to metric-frame existence}
\author{Run \texttt{fa\_banach\_001}}
\date{2026-06-28}

\begin{document}
\maketitle

\section*{Source and Target}

The source paper is K. Mahesh Krishna's arXiv:2202.13697 paper
\cite{Krishna2022}.  Its title is
\emph{Metric, Schauder and Operator-Valued Frames}.  In Section 2.1, after
defining metric frames, the source
states on PDF page 68 that, after the definition of metric frame, ``the first
question which comes is the existence.''  It then explains that the analogous
Banach-frame existence theorem is known for separable Banach spaces, while the
metric-space setting is not known in that form.

\begin{center}
\includegraphics[width=\textwidth,height=0.78\textheight,keepaspectratio]{figures/open_problem_crop.png}
\end{center}

\section*{Definitions}

We use the source's Definition 2.1.10.  A metric frame for a metric space
\(\mathcal M\), with respect to a BK-space \(\mathcal M_d\), consists of
Lipschitz functions \(f_n:\mathcal M\to\mathbb K\) and a Lipschitz
reconstruction operator \(S:\mathcal M_d\to\mathcal M\) such that the analysis
map
\[
  \theta_f(x)=(f_n(x))_{n\ge 1}\in\mathcal M_d
\]
is bi-Lipschitz onto its image and
\[
  S(\theta_f(x))=x,\qquad x\in\mathcal M.
\]
A BK-space is, in particular, a normed linear space and hence is path
connected.

\section*{Proof Intuition}

The reconstruction operator in a metric frame is stronger than a formal
coordinate inverse: it is a Lipschitz map from the whole Banach sequence space
\(\mathcal M_d\) onto \(\mathcal M\).  Since a Banach space is linearly path
connected, composing the line segment between two analysis vectors with
\(S\) gives a path in \(\mathcal M\).  Thus metric-frame reconstruction forces
the underlying metric space to be path connected.  A two-point metric space is
separable but not path connected, so it cannot admit such a frame.

\begin{theorem}
If a metric space \(\mathcal M\) admits a metric frame in the sense of
Definition 2.1.10 of \cite{Krishna2022}, then \(\mathcal M\) is Lipschitz path
connected.
\end{theorem}

\begin{proof}
Let \((\{f_n\}_{n\ge 1},S)\) be a metric frame for \(\mathcal M\) with
coefficient BK-space \(\mathcal M_d\).  Fix \(x,y\in\mathcal M\).  Put
\(u=\theta_f(x)\) and \(v=\theta_f(y)\).  Since \(\mathcal M_d\) is a normed
linear space, the segment
\[
  \sigma(t)=(1-t)u+tv,\qquad 0\le t\le 1,
\]
is a Lipschitz path in \(\mathcal M_d\).  The composition
\(\gamma=S\circ\sigma:[0,1]\to\mathcal M\) is Lipschitz because \(S\) is
Lipschitz.  Moreover,
\[
  \gamma(0)=S(\theta_f(x))=x,\qquad
  \gamma(1)=S(\theta_f(y))=y.
\]
Thus every two points of \(\mathcal M\) are joined by a Lipschitz path.
\end{proof}

\begin{corollary}
The two-point metric space \(\mathcal M=\{0,1\}\), with
\(d(0,1)=1\), is a separable metric space which admits no metric frame in the
sense of Definition 2.1.10 of \cite{Krishna2022}.
\end{corollary}

\begin{proof}
The space \(\{0,1\}\) is separable.  It is not path connected: the continuous
image of the connected interval \([0,1]\) into the discrete two-point metric
space must be a singleton.  This contradicts the theorem if a metric frame
existed.
\end{proof}

\section*{Verification Notes}

The argument uses only the Lipschitz reconstruction clause in the definition
of metric frame.  It does not use the lower or upper frame inequalities except
to ensure that the analysis vectors lie in the BK-space.

This counterexample does not contradict the source's Theorem 2.1.17.  That
theorem gives Lipschitz coordinate functions and a reconstruction function, but
its statement does not require the reconstruction function to be Lipschitz.
The obstruction above applies precisely to the full metric-frame definition,
where \(S:\mathcal M_d\to\mathcal M\) is required to be Lipschitz on the whole
BK-space.

No computational checks were needed.  The source PDF and evidence crop are
included in this packet.

\section*{Novelty and Literature Check}

The run indexes were searched before promotion for arXiv:2202.13697 and for
core phrases including \texttt{metric frame}, \texttt{Lipschitz retract}, and
\texttt{path connected}.  No prior packet or attempt for this target was
found.  Web and arXiv searches on 2026-06-28 for the exact source title and
phrases such as \texttt{metric frame separable metric}, \texttt{metric frame
Lipschitz retract}, and \texttt{metric frame path connected} found the source
paper, the related arXiv:2011.01870 paper, and adjacent Lipschitz-frame work,
but no duplicate of this path-connectedness obstruction.

\section*{Human Review Recommendation}

Check the source definition of metric frame, especially that the reconstruction
operator \(S\) is required to be Lipschitz on all of \(\mathcal M_d\).  Under
that definition, the proof above is immediate and the two-point space refutes
the broad existence question for arbitrary separable metric spaces.

\begin{thebibliography}{9}

\bibitem{Krishna2022}
K. Mahesh Krishna.
\newblock \emph{Metric, Schauder and Operator-Valued Frames}.
\newblock arXiv:2202.13697, 2022.

\end{thebibliography}

\end{document}
